\documentclass[12pt,a4paper,psamsfonts]{amsart}
\usepackage{amssymb,mathtools}
\usepackage{calc}
\setlength{\topmargin}{0cm}
\setlength{\oddsidemargin}{0cm}
\setlength{\evensidemargin}{0cm}
\setlength{\marginparwidth}{0cm}
\setlength{\marginparsep}{0cm}
\setlength{\textheight}{\paperheight - 2in - 35pt}
\setlength{\textwidth}{\paperwidth - 2in}
\setlength{\headheight}{12.5pt}
\setlength{\headsep}{25pt}
\setlength{\footskip}{30pt}
\renewcommand{\baselinestretch}{1.3}
\pagestyle{headings}
\raggedbottom
\setlength{\emergencystretch}{2em}
\usepackage[colorlinks=true,urlcolor=blue,linkcolor=blue,citecolor=blue]{hyperref}
\newtheorem*{theorem}{Theorem}
\theoremstyle{definition}
\newtheorem*{setup}{Setup}
\theoremstyle{plain}
\newtheorem{lemma}{Lemma}[section]
\newtheorem{proposition}[lemma]{Proposition}
\newtheorem{corollary}[lemma]{Corollary}
\newcommand{\Q}{\mathbb Q}
\newcommand{\C}{\mathbb C}
\newcommand{\OO}{\mathcal O}
\newcommand{\mult}{\operatorname{mult}}
\newcommand{\nef}{\operatorname{n}}

\usepackage{tikz-cd}
\newcommand{\R}{\mathbb R}
\newcommand{\Z}{\mathbb Z}
\newcommand{\PP}{\mathbb P}
\DeclareMathOperator{\Pic}{Pic}
\DeclareMathOperator{\Supp}{Supp}
\DeclareMathOperator{\rk}{rk}
\DeclareMathOperator{\ord}{ord}
\DeclareMathOperator{\vol}{vol}
\DeclareMathOperator{\Sym}{Sym}
\DeclareMathOperator{\Bir}{Bir}
\DeclareMathOperator{\codim}{codim}
\def\UrlBreaks{\do\/\do-\do.}
\title{Nef Abundance: Framework -- Jet Separation in General Dimension}
\date{Temporary draft -- October 11, 2026}
\begin{document}
\maketitle

\section{Jet separation}

Fix an integer $d\ge2$. We use the following source inputs,
denoted collectively by $(\mathsf I_d)$.
\begin{itemize}
\setlength{\itemsep}{2pt}
\setlength{\parsep}{0pt}
\item For every projective klt $\Q$-pair $(Y,B)$ of dimension
$j<d$ and every nef $\mathbb{Q}$-Cartier divisor $M$, if $K_Y+B$ is
pseudo-effective and, for some $a\in\Q_{>0}$,
\[
 (K_Y+B+aM)\cdot C>0
\]
for every integral curve through a sufficiently very general point,
then $K_Y+B+cM$ is big for every $c\in\Q_{>0}$.
\item Every projective log canonical $\Q$-pair of dimension at most
$d$ with nef adjoint has semiample adjoint.
\item Every projective $\Q$-factorial generalized klt $\Q$-pair of
dimension at most $d$, with fixed nef $\mathbb{Q}$-Cartier b-divisor
data and pseudo-effective adjoint, admits a terminating choice of
minimal-model program with a nef endpoint. The program extracts no
divisors, preserves $\Q$-factorial generalized klt models, and on a
common resolution the initial adjoint equals the final adjoint plus
an effective divisor exceptional over the final model.
\end{itemize}
The first input is the lower-dimensional assertion $(P_j)$ of
\cite[Section~3.1]{NS}. The last two are the companion inputs
used in \cite[Section~2]{NS}; here they are assumptions, not
independently established higher-dimensional theorems.

Let $X$ be a smooth projective variety over $\C$ of dimension $d$,
with $K_X\sim0$ and $q(X)=0$. Let $D$ be a nef Cartier divisor
such that $\nef(X,D)=d$ and $D^d=0$. Assume $(\mathsf I_d)$.

\begin{setup}
\leavevmode
\begin{itemize}
\setlength{\itemsep}{2pt}
\setlength{\parsep}{0pt}
\item Fix an ample Cartier divisor $A$.
\item Fix $\epsilon\in\Q_{>0}$ sufficiently small. For integers
$n\ge1$, put
\[
 a_n=((D+A/n)^d)^{1/d},\qquad
 r_n=\max\{1,\lfloor\epsilon/a_n+1/2\rfloor\},\qquad
 P_n=r_n(D+A/n).
\]
Then $P_n$ is ample and $\mathbb{Q}$-Cartier, $r_n\to\infty$, and
$P_n^d\to\epsilon^d$. Indeed, $a_n\to0$ and
$|r_na_n-\epsilon|\le a_n/2$ for all sufficiently large $n$.
The symbol $n$ indexes the sequence; the dimension is $d$.
\item Let $p_i:X\times X\to X$ be the projections. Put
\[
 Z=\PP_{X\times X}
       \bigl(p_1^*\OO_X(D)\oplus p_2^*\OO_X(D)\bigr),
 \qquad \pi:Z\to X\times X.
\]
Let $S_1,S_2$ be the two summand sections and set
$U=Z\setminus(S_1\cup S_2)$. Thus $\dim Z=2d+1$.
\item Set $\xi=c_1(\OO_Z(1))$. For an integer $s>0$, put
\[
 R_{n,s}=(p_1\circ\pi)^*P_n+(p_2\circ\pi)^*P_n+s\xi.
\]
We use $\pi_*\OO_Z(a)=\operatorname{Sym}^a
(p_1^*\OO_X(D)\oplus p_2^*\OO_X(D))$ for $a\ge0$.
Whenever $|kR_{n,s}|$ is used, $k>0$ is an integer with
$kP_n$ Cartier.
\item Choose $\epsilon$ sufficiently small that
\[
 (14\epsilon)^d<\tfrac14,\qquad40\epsilon<\tfrac38.
\]
The integer $s$ is chosen after $\epsilon$ and before $n$ tends
to infinity. Any jet-separating multiple is chosen after $n$.
\end{itemize}
\end{setup}

\medskip
\noindent\textbf{Jet separation.}
Let $W$ be smooth, let $L$ be a Cartier divisor, and let $z\in W$.
For $j\ge0$, $|L|$ separates $j$-jets at $z$ if
\[
 \forall u\in\OO_{W,z}(L)\quad
 \exists\sigma\in H^0(W,\OO_W(L))\quad
 \sigma_z-u\in\mathfrak m_z^{j+1}\OO_{W,z}(L).
\]

\begin{lemma}\label{lem:general-jets}
There exists a positive integer $s$ such that, for every $n\gg0$,
$|kR_{n,s}|$ separates $(2d+2)k$-jets at a very general point
$z\in U$, for all sufficiently divisible positive integers $k$.
\end{lemma}

\medskip
\noindent\textbf{Draft status.}
This is a conditional, source-based proof framework following
\cite[Sections~4--6]{NS}. The source inputs $(\mathsf I_d)$ and
the inherited moving-centre and Frobenius arguments have not been
independently verified here. No unconditional higher-dimensional
nonvanishing or abundance theorem is asserted.
\section{Auxiliary estimates}
\label{sec:jet-tools}

The geometric argument will produce subvarieties with bounded degree and
bounded canonical intersection.  We first explain how these bounds
contradict a growing lower bound for the degrees of curves.  We then
record the moving-center results from \cite{LA} that supply the required
subvarieties.  All varieties in this section are over $\C$.

\subsection{The bounded-degree obstruction}

\begin{lemma}
\label{lemma:bounded-curves}
For every integer $f\geq1$ there is an integer $m_f>0$ with the
following property.  Let $Y$ be a smooth integral projective variety
of dimension $f$, let $D_0$ be a reduced simple normal crossing divisor
or the zero divisor, and let $L_0$ be a nef integral Cartier divisor
of full nef dimension on $Y$.  Suppose that $P_0$ is a nef $\mathbb{Q}$-Cartier divisor and that $r,C_0,C_1$ are positive real numbers satisfying
\[
 0<P_0^f\leq C_0,\qquad P_0-rL_0\text{ is nef}.
\]
If $A_0=K_Y+D_0+L_0$ is big and
\begin{equation}
\label{j:adjoint-ratio}
 A_0\cdot P_0^{f-1}\leq C_1P_0^f,
\end{equation}
then
\begin{equation}
\label{j:bounded-curve-obstruction}
 r\leq(m_fC_1)^{f-1}C_0.
\end{equation}
The integer $m_f$ depends only on $f$, and not on the variety or
the divisors.  For dimensions $1\leq f\leq d$, one may replace all
$m_f$ by a single integer $m$ depending only on $d$.
\end{lemma}

\begin{proof}
The pair $(Y,D_0)$ is log canonical, its boundary coefficients belong
to $\{0,1\}$, and its nef polarizing divisor $L_0$ has Cartier index
one.  By \cite[Theorem~1.3]{BZ}, an integer $m_f$ depending only on
these data makes $|m_fA_0|$ birational.  Here $A_0$ is integral
Cartier, so no rounding or additional denominator is involved.

Resolve the base ideal by a projective birational morphism
$\mu:\widetilde Y\to Y$, with $\widetilde Y$ smooth, and write
\[
 \mu^*(m_fA_0)=H+F,
\]
where $H$ is the basepoint-free moving part and $F$ is effective.
The morphism defined by $H$ is birational onto its image; in
particular $H$ is nef and big.  In the rest of this proof, $P_0$
and $L_0$ denote their pullbacks to $\widetilde Y$.  Effectivity
of $F$ and nefness of $P_0$ give
\begin{equation}
\label{j:moving-adjoint-degree}
 H\cdot P_0^{f-1}\leq m_fC_1P_0^f.
\end{equation}

For $f\geq2$, put $s_i=P_0^{f-i}\cdot H^i$ for $0\leq i\leq f$.
Both divisors are nef and big, so the $s_i$ are positive.  The
Khovanskii--Teissier inequalities
\cite[Corollary~1.6.3(i) and Example~1.6.4]{LazarsfeldPAGI}
$s_i^2\geq s_{i-1}s_{i+1}$ show
that the ratios $s_i/s_{i-1}$ decrease.  Hence
\begin{equation}
\label{j:mixed-degree}
 P_0\cdot H^{f-1}
 \leq P_0^f\left(\frac{H\cdot P_0^{f-1}}{P_0^f}\right)^{f-1}
 \leq(m_fC_1)^{f-1}C_0.
\end{equation}
Choose a very general point $y\in\widetilde Y$ over the
isomorphism opens of both $\mu$ and the morphism defined by $H$.
We also choose $\mu(y)$ in the very general locus
that tests full nef dimension of $L_0$.  Pull back $f-1$ general
hyperplanes through the image of $y$ under the morphism defined
by $H$.  This subsystem has only the base point $y$, so the
general members meet properly and define an effective curve cycle
of class $H^{f-1}$; the component $C$ through $y$ is a curve
because their local equations are independent at $y$.
Nefness gives
\[
 P_0\cdot C\leq P_0\cdot H^{f-1}.
\]
The curve $C$ is not contracted by $\mu$, and its image passes
through the chosen very general point.  Full nef dimension and
integrality therefore give $L_0\cdot C\geq1$.  Since
$P_0-rL_0$ is nef, we obtain $P_0\cdot C\geq r$, which proves
\eqref{j:bounded-curve-obstruction}.

If $f=1$, take $C=Y$ directly.  Then
$P_0\cdot Y\geq rL_0\cdot Y\geq r$, and the required bound is
$r\leq C_0$.  Finally, \cite[Theorem~1.3]{BZ} applies to every
multiple of its prescribed integer.  A common multiple of the
finitely many integers $m_1,\ldots,m_d$ therefore gives the last
assertion.
\end{proof}

In particular, the hypotheses of Lemma~\ref{lemma:bounded-curves}
cannot hold along a sequence with $r\to\infty$ and with $C_0,C_1$
bounded, even if the varieties and divisors vary.  The following
observation explains how full nef dimension and pseudo-effectivity
of the canonical divisor pass to the moving subvarieties used below.

\begin{lemma}
\label{j:very-general-images}
Let $X$ be a normal integral projective variety that is not uniruled,
and let $L$ be a nef Cartier divisor of full nef dimension on $X$.
There is a complement $X^\circ$ of a countable union of proper
closed subsets of $X$ with the following property.  If $Y$ is a
smooth integral projective variety and $g:Y\to X$ is generically
finite onto a positive-dimensional image meeting $X^\circ$, then
$g^*L$ has full nef dimension, $Y$ is not uniruled, and $K_Y$ is
pseudo-effective.
\end{lemma}

\begin{proof}
Choose $X^\circ$ so that every curve through any of its points has
positive $L$-degree and no rational curve meets $X^\circ$.  The first
condition is the very general characterization of full nef
dimension.  For the second, rational curves are parametrized by
countably many algebraic families; the closure of the evaluation
image of each such family is proper because $X$ is not uniruled.

Since $g(Y)$ meets $X^\circ$, it is contained in none of the closed
sets excluded in its definition.  A very general point of $Y$ thus
maps into $X^\circ$ and lies in the quasi-finite locus of $g$.
Every curve through such a point has a curve as its image, so its
$g^*L$-degree is positive.  This proves full nef dimension.  A
covering family of rational curves on $Y$ would similarly give a
nonconstant rational curve meeting $X^\circ$, a contradiction.
The pseudo-effectivity of $K_Y$ now follows from \cite{BDPP}.
\end{proof}

\subsection{Subadjunction and high-multiplicity components}

We use the following five results from the section
\emph{Tools for moving base components} of \cite{LA}.  Their
hypotheses distinguish generic information along a center from
global information about the ambient variety.  Intersections on
an integral subvariety are computed after pullback to its
normalization and a smooth projective resolution.  A base ideal
is \emph{isolated at the generic point} of a subvariety when its
localization there is primary for the maximal ideal.

\begin{lemma}
\label{j:subadjunction}
Let $Z$ be a projective $\Q$-factorial klt variety, let
$\Theta\geq0$ be a rational divisor, and let $V\subset Z$ be
integral of dimension $f>0$.  Suppose that $(Z,\Theta)$ is log
canonical at the generic point of $V$ and has a divisor of log
discrepancy zero with center $V$.  If $P$ is a nef $\mathbb{Q}$-Cartier
divisor on $V$ and $\rho:V^*\to V$ is a smooth projective resolution
factoring through its normalization, then
\[
 K_{V^*}\cdot(\rho^*P)^{f-1}
 \leq (K_Z+\Theta)|_V\cdot P^{f-1}.
\]
\end{lemma}

Only generic log canonicity is required.  The proof in \cite{LA}
uses the dlt modification of \cite[Theorem~2.10]{FH}, whose
truncated boundary leaves an effective surplus supported over
the non-lc locus.  On a minimal stratum over $V$, adjunction
gives a pair that is klt over the generic point of its
connected-fiber base.  This suffices for the klt-trivial
canonical bundle formula and its b-nef moduli divisor
\cite[Definition~3.1 and Theorem~3.5]{FG}.  The effective
discriminant, the moduli divisor, and finite ramification have
nonnegative intersections with the pulled-back nef class.
The comparison uses canonical Weil cycles, so it requires no
$\Q$-Gorenstein assumption on the normalization of $V$.

\begin{lemma}
\label{j:base-component}
Let $Z$ be a projective $\Q$-factorial klt variety of dimension
$e$, let $R$ be a nef $\mathbb{Q}$-Cartier divisor, and let $k>0$
be an integer such that $kR$ is Cartier.  Let
$0\ne\mathcal H\subset H^0(Z,\OO_Z(kR))$ be a linear subspace
with proper base locus and base ideal $\mathfrak b$.  Suppose that
$V$ is an integral base-locus component of dimension $f<e$,
isolated at its generic point $\eta_V$, and that
$\eta_V\in Z_{\mathrm{sm}}$.  Let
$\mathfrak m$ be the maximal ideal of $\OO_{Z,\eta_V}$.
If $\mathfrak b\OO_{Z,\eta_V}\subset\mathfrak m^h$ for an
integer $h>0$, there are a rational number $0<c\leq(e-f)k/h$ and an
effective rational divisor $\Theta\sim_{\Q}cR$ such that
$(Z,\Theta)$ is log canonical at $\eta_V$ and has a log canonical
place centered on $V$.  The degree and canonical bounds are
\begin{equation}
\label{j:base-bounds}
 \begin{split}
 R^f\cdot V&\leq(k/h)^{e-f}R^e,\\
 K_{V^*}\cdot R^{f-1}
   &\leq(K_Z+cR)|_V\cdot R^{f-1}\qquad(f>0),
 \end{split}
\end{equation}
where $V^*$ is a smooth projective resolution.  The degree bound
also holds for $f=0$.
\end{lemma}

The degree estimate is the local multiplicity bound $h^{e-f}$
combined with successive general cuts by the given linear series.
The boundary is obtained by averaging general members with total
coefficient equal to the generic log canonical threshold of
$\mathfrak b$.  Thus it is an actual divisor to which
Lemma~\ref{j:subadjunction} applies.

\subsection{Full jet kernels and restriction to a fiber}

\begin{lemma}
\label{j:moving-kernel}
Let $Z$ be an integral projective variety of dimension $e$, let
$R$ be a nef, big, semiample $\mathbb{Q}$-Cartier divisor, and fix
an integer $k>0$ such that $kR$ is Cartier.  Choose positive real
numbers $\tau_0<\cdots<\tau_e$ with consecutive gap $\delta>0$.
For $x\in Z_{\mathrm{sm}}$, put
\[
 \mathcal H_i(x)=H^0\bigl(Z,\OO_Z(kR)
                  \otimes\mathfrak m_x^{\lceil k\tau_i\rceil}\bigr).
\]
Suppose that for very general $x$ all these spaces are nonzero and
the base locus of $\mathcal H_0(x)$ has a positive-dimensional
component through $x$.  If $k\delta\geq4$, then, after an
algebraic parameter extension and restriction to nonempty opens,
there are an integral parameter variety $T$ with a dominant
marked-point map $x:T\to Z_{\mathrm{sm}}$, a fixed index
$i\in\{0,\ldots,e-1\}$, and a family of integral subvarieties
$V_t$ containing $x(t)$ and common to the base loci of
$\mathcal H_i(x(t))$ and $\mathcal H_{i+1}(x(t))$.  In particular,
the incidence
\[
 \Gamma=\{(t,z):z\in V_t\}\subset T\times Z
\]
dominates $Z$.

For general $t$, the dimension $f$ of $V_t$ lies in $[1,e-1]$,
$V_t$ is an isolated component of the higher base locus at its
generic point, and the entire higher-kernel base ideal is
contained there in the ordinary power $\mathcal I_{V_t}^h$, where
\[
 h=\lceil k\delta/2\rceil,\qquad k/h\leq2/\delta.
\]
The incidence base ideal is likewise contained in
$\mathcal I_\Gamma^h$ on a dense open of $\Gamma$.
These containments are taken in the smooth ambient open.
If $Z$ is also $\Q$-factorial and klt, then
\begin{equation}
\label{j:moving-bounds}
 \begin{split}
 R^f\cdot V_t&\leq(2/\delta)^{e-f}R^e,\\
 K_{V_t^*}\cdot R^{f-1}
   &\leq(K_Z+cR)|_{V_t}\cdot R^{f-1},
       \qquad 0<c\leq2(e-f)/\delta.
 \end{split}
\end{equation}
The top intersection on $V_t$ is positive when $V_t$ meets the
open on which the big semiample contraction is an isomorphism.
\end{lemma}

The use of the full kernels matters: differentiation in the
parameter lowers vanishing at the moving mark and still gives a
global section in the lower kernel.  Incidence dominance makes
these parameter directions span the normal directions.  The
vanishing of all derivatives of order less than $h$ then gives
the ordinary-power containment, which can be restricted to
general fibers as follows.

\begin{lemma}
\label{j:fiber-restriction}
In the setting of Lemma~\ref{j:moving-kernel}, let $g:Z\to B$ be
a morphism.  There is a dense open of the incidence with the
following property.  Choose $(t,z)$ in this open, put $b=g(z)$,
and suppose that $V_t$ is smooth over a smooth open of its actual
image at $z$.  Let $F$ be the reduced component through $z$ of
the scheme fiber of $V_t\to g(V_t)$.  Near $z$ in $Z_b$, the
restricted higher-kernel base ideal has support exactly $F$ and
is contained in the ordinary power $\mathcal I_F^h$.

If $F$ is proper in an integral ambient fiber component containing
it, the restricted series is nonzero there and $F$ is an isolated
base component at its generic point.  Applying
Lemma~\ref{j:base-component} on that ambient component additionally
requires its projectivity, $\Q$-factoriality, klt singularities,
and smoothness at the generic point of $F$.
\end{lemma}

Here the fiber is chosen through a general incidence point.  The
statement does not assert the same conclusion for an arbitrary
special fiber.  Smoothness over the actual image identifies the
restricted ideal of $V_t$ with the reduced ideal of $F$ locally;
this is what preserves ordinary powers.

The $\Q$-factoriality needed for the base-component estimates can
also be supplied by a small projective $\Q$-factorialization
$\sigma:Z'\to Z$ of a klt ambient variety.  This passage is made
only for components that meet the isomorphism open of $\sigma$.
Their strict transforms retain the generic base-ideal and
multiplicity conditions for the pulled-back series, and the
testing class $\sigma^*R$ is nef.  Since
$K_{Z'}=\sigma^*K_Z$, projection formula gives the same degree
and canonical-intersection bounds on a common resolution.
The pullback testing class need not be ample; nefness is the
property used in these estimates.  Passing to higher resolutions
does not change the canonical intersections, because their
additional exceptional divisors have zero intersection with
the pulled-back testing class to the required power.

\subsection{Finite covers with a fixed branch complement}

A bound on the degree of a branched cover does not make the
possible covers finite.  The final companion input supplies the
needed fixed branch complement from a family whose branch
divisors do not sweep the base.

\begin{lemma}
\label{j:finite-covers}
Let $Y$ be a normal integral projective variety, let
$\mathcal V\to T$ be a smooth projective family with integral
fibers over an integral parameter variety, and let
$g:\mathcal V\to Y$ be a morphism.  Suppose that the general
fiber maps $g_t:V_t\to Y$ are dominant and generically finite of
degree at most a fixed integer $e$.

After a finite parameter extension and shrinking, spread the
geometric generic components of the zero divisor of the relative
Jacobian over $Y_{\mathrm{sm}}$ to divisors
$\mathcal R_1,\ldots,\mathcal R_N$ in $\mathcal V$.
Retain this notation after the base change.  Let $I$ consist of
the indices $j$ for which
$\overline{g_t((\mathcal R_j)_t)}$ is a divisor in $Y$ for general
$t$.  Assume that
\[
 \overline{g(\mathcal R_j)}\ne Y\qquad(j\in I).
\]
After shrinking again, there is a smooth nonempty open
$Y^\circ\subset Y$ over which every general finite normal Stein
cover of $g_t$ is finite \'etale.  Only finitely many such normal
covers of $Y$ occur, up to isomorphism over $Y$.
The conclusion holds simultaneously for finitely many
projections; the open and the finite lists may depend on all
fixed parameters of the family.
\end{lemma}

To see the role of the hypothesis, remove $Y_{\mathrm{sing}}$
and the finitely many proper image closures
$\overline{g(\mathcal R_j)}$ for $j\in I$.  Every branch prime
of a general Stein cover is accounted for by one of these
relative Jacobian components.  Purity therefore makes the
restricted covers \'etale.  The fundamental group of the
resulting smooth complex variety is finitely generated, so it
has only finitely many covers of degree at most $e$.  Finally,
a finite normal cover of $Y$ is determined by its function
field, hence by its restriction to this fixed open.

\begin{lemma}\label{prop:full-positive}
Assume $(\mathsf I_d)$. Let $(Y,B)$ be a projective klt $\Q$-pair
of dimension at most $d$, with $T=K_Y+B$ pseudo-effective and
$\kappa(Y,T)>0$. Let $M$ be nef and $\mathbb{Q}$-Cartier. If
$(T+aM)\cdot C>0$ for every integral curve through a sufficiently
very general point, for some $a>0$, then $T+cM$ is big for every
rational $c>0$.
\end{lemma}
\begin{proof}
Apply the source reduction for positive Iitaka dimension
\cite[Proposition~3.3]{NS}. Its model-program and ordinary abundance
inputs are included in $(\mathsf I_d)$; its fibre dimensions are
strictly smaller than $d$. This reduction is retained as a
source-dependent input here.
\end{proof}

\section{Geometric estimates}\label{sec:geometry}
In the technical arguments below, put $L=D$, $t=1/n$, $r=r_n$,
$P=P_n$, and $q=s$. The letter $L$ keeps the source notation for
the fixed nef Cartier divisor; $q$ is the bundle parameter.
\begin{equation}\label{eq:scaling}
 P=r(L+tA),\qquad r\to\infty,\qquad P^d\to\epsilon^d.
\end{equation}
For each fixed $q$, discard finitely many terms so that
\begin{equation}\label{j:volume-window}
 \epsilon^d/2\le P^d\le2\epsilon^d,\qquad r>q+1.
\end{equation}
Every curve through a very general point of $X$ has positive integral
$L$-degree, hence $P$-degree at least $r$. Thus a covering family of
curves of bounded $P$-degree will give a contradiction.

We specify how the same obstruction applies to subvarieties and
covers. A resolution $W\to X$ has an effective canonical divisor:
pull back a nowhere-zero canonical form on the smooth locus of $X$
and use terminality. Hence $W$, and therefore $X$, is not uniruled
\cite{BDPP}. By Lemma~\ref{j:very-general-images}, if a smooth
projective variety $Y$ maps generically finitely to a
positive-dimensional subvariety containing a sufficiently very general
point of $X$, then $Y$ is not uniruled and the pulled-back Cartier
divisor $L_Y$ has full nef dimension. In particular, $K_Y$ is
pseudo-effective. If $f=\dim Y<d$, the induction hypothesis $(P_f)$
applied to $(Y,0)$ gives
\begin{equation}\label{j:lower-big}
                         K_Y+L_Y\ \text{big}.
\end{equation}
All occurrences of ``very general'' below include the fixed
conditions of Lemma~\ref{j:very-general-images}.

Consequently, for $f<d$ it suffices to find on such a $Y$ a nef
testing divisor $P_Y$ with
\[
 0<P_Y^f\le C_0,\qquad P_Y-rL_Y\ \text{nef},\qquad
 K_YP_Y^{f-1}\le C P_Y^f,
\]
where $C_0,C$ are independent of $t$. Indeed
$L_YP_Y^{f-1}\le P_Y^f/r$, so Lemma~\ref{lemma:bounded-curves}
applies. We will also use that lemma in dimension $d$, after proving
the required bigness separately. For the rest of this section, a
bound is called uniform if it is independent of $t$; it may depend
on parameters fixed before $t$ tends to zero. Since the dimensions
lie in a finite range, Lemma~\ref{lemma:bounded-curves} excludes all
such configurations once $r$ is sufficiently large.

\subsection{The projective bundle and its volume}

We now seek a jet system whose order is much larger than the scalar
bound. The additional direction comes from the fibres of a
projective bundle; its two summands keep track of the two copies of
$X$. For a divisor on $X$, subscripts $1,2$ denote pullback from the
two factors of $X\times X$. Set
\begin{equation}\label{eq:bundle}
 Z=\PP\bigl(\OO_X(L)_1\oplus\OO_X(L)_2\bigr)
       \xrightarrow{\pi}X\times X,
 \qquad \xi=c_1(\OO_Z(1)),
 \qquad R=P_1+P_2+q\xi,
\end{equation}
where $q$ is a positive integer to be fixed shortly. We use the
convention
$\pi_*\OO_Z(a)=\operatorname{Sym}^a(\OO_X(L)_1\oplus\OO_X(L)_2)$
for $a\ge0$. The tautological class $\xi$ is nef, and $R$ is ample.
The product and projective-bundle descriptions show that $Z$ is klt.
The two natural sections, called the axes, have classes
$\xi-L_1$ and $\xi-L_2$. The canonical bundle formula gives
\begin{equation}\label{j:axes-canonical}
                         K_Z\sim-\text{(sum of the two axes)}.
\end{equation}
We call the complement of the axes the torus open.

The fibre of either projection $Z\to X$ over a smooth point is,
up to a constant one-dimensional factor,
\[
 Z_{\mathrm{sl}}=\PP(\OO_X\oplus\OO_X(L)),
 \qquad R_{\mathrm{sl}}=P+q\xi.
\]
Its axes have classes $\xi$ and $\xi-L$, and its canonical divisor
is again minus their sum. The slice is projective and klt. In
particular it has the singularity hypotheses needed to apply the
base-component estimates, using a small $\Q$-factorialization if
necessary as explained after Lemma~\ref{j:moving-kernel}.

Write $e=2d+1$. There are positive constants depending only on $d$
such that, for each fixed $q$ and all sufficiently small $t$,
\begin{equation}\label{j:bundle-volumes}
 c_nq\epsilon^{2d}\le R^e\le C_nq\epsilon^{2d},
 \qquad R_{\mathrm{sl}}^{d+1}\le C_{\mathrm{sl},d}q\epsilon^d.
\end{equation}
To see this, expand in nef classes and use
\[
 \pi_*(\xi^j)=\sum_{a=0}^{j-1}L_1^aL_2^{j-1-a}
 \quad(j\ge1).
\]
The term containing exactly one $\xi$ is
$(2d+1)\binom{2d}{d}q(P^d)^2$, giving the lower bound.
For the upper bound, every occurrence of $L_i$ can be bounded in
mixed nef intersections by $P_i/r$. After factoring out $q$, the
remaining powers of $q/r$ are bounded by one once $r\ge q$.
There are only finitely many terms depending on $d$. The slice
calculation is the same, using $\pi_*(\xi^j)=L^{j-1}$.
Together with \eqref{j:volume-window}, this proves
\eqref{j:bundle-volumes}.

For a nef $\mathbb{Q}$-Cartier divisor $R$ and a smooth point $z$, its
\emph{Seshadri constant} is
\[
 \varepsilon(R;z)=\inf_{C\ni z}\frac{R\cdot C}{\mult_z C},
\]
where $C$ ranges over integral curves through $z$. Thus a lower bound
for this constant controls degree relative to multiplicity at the point.

\begin{proposition}\label{prop:large-seshadri}
For a suitable fixed positive integer $q$, and all sufficiently
small $t$ in \eqref{eq:scaling}, there is a smooth point $z\in Z$
in the torus open, lying over smooth points of both factors, such
that
\begin{equation}\label{j:large-seshadri}
                         \varepsilon(R;z)>2d+2.
\end{equation}
More generally, for any prescribed dense open subsets $U_1,U_2$
of the smooth locus of $X$, the point may be chosen over
$U_1\times U_2$. In particular, these opens may be the
isomorphism loci of fixed resolutions of $X$.
\end{proposition}

We prove the proposition in the remainder of this section. Choose a
positive rational gap $\delta$ such that
\begin{equation}\label{j:slice-gap}
             (2/\delta)^nC_{\mathrm{sl},d}\epsilon^d<1.
\end{equation}
Next choose $q$ so large that $e+1$ increasing levels of gap $\delta$,
with lowest level strictly greater than $2d+2$, lie strictly below
$(c_nq\epsilon^{2d})^{1/e}$. These choices precede the limit
$t\to0$.

Fix the prescribed opens $U_1,U_2$, if any. Suppose that
\eqref{j:large-seshadri} fails throughout their smooth torus
preimage for arbitrarily small $t$. For any such $t$, asymptotic
Riemann--Roch and the elementary
jet count at a smooth point show that all the full kernels at the
chosen levels are nonzero for sufficiently large divisible $k$.
Indeed the highest level has $e$-th power smaller than $R^e$.
Failure of \eqref{j:large-seshadri} supplies, through each very
general point, a curve with degree divided by multiplicity smaller
than the lowest level. Every section in the lowest kernel vanishes
identically on that curve: otherwise its intersection with the
curve is at least its vanishing order times the curve's
multiplicity, contradicting the degree inequality. Thus the lowest
base locus has a positive-dimensional component through the mark.
Take $k\delta\ge4$ and apply Lemma~\ref{j:moving-kernel}.

We obtain a family of proper integral components $V\subset Z$ whose
incidence dominates $Z$. For its general members, of dimension $f$,
the estimates are
\begin{equation}\label{j:total-component-bounds}
 \begin{split}
 0<R^f\cdot V&\le(2/\delta)^{e-f}R^e,\\
 K_{V^*}R^{f-1}&\le(K_Z+cR)|_V R^{f-1},
 \qquad 0<c\le2(e-f)/\delta.
 \end{split}
\end{equation}
Here and below classes on a resolution are understood by pullback.
The positivity follows from ampleness of $R$. General members meet
the torus open, so \eqref{j:axes-canonical} makes the ambient
canonical contribution nonpositive. All constants in these bounds
are uniform now that $q,\epsilon,\delta$ have been fixed.

\subsection{Excluding positive-dimensional slice fibres}

We first constrain the two projections of $V$ to $X$. For either
projection, consider the component $F$ of a general fibre of
$V\to\operatorname{im}(V\to X)$ through a general incidence point.
Choose that incidence point in the opens of
Lemma~\ref{j:fiber-restriction}, over very general points of both
factors of $X$, and in the smooth torus open. We claim that
\begin{equation}\label{j:no-intermediate-fibre}
                         \dim F\notin\{1,\ldots,d\}.
\end{equation}

Suppose instead that $1\le f_{\mathrm{sl}}=\dim F\le d$.
The ambient fibre is $Z_{\mathrm{sl}}$, of dimension $d+1$.
The restriction lemma gives a nonzero restricted subseries with
$F$ an isolated base component and base ideal contained generically
in $\mathcal I_F^h$, where $h=\lceil k\delta/2\rceil$.
Generic smoothness is imposed over the actual image of $V$; this is
why a general fibre through the incidence point was chosen.
The point is also smooth in the ambient slice. Applying
Lemma~\ref{j:base-component} on that slice gives
\begin{equation}\label{j:slice-bounds}
 \begin{split}
 0<R_{\mathrm{sl}}^{f_{\mathrm{sl}}}\cdot F
   &\le(2/\delta)^{d+1-f_{\mathrm{sl}}}
                      R_{\mathrm{sl}}^{d+1},\\
 K_{F^*}R_{\mathrm{sl}}^{f_{\mathrm{sl}}-1}
   &\le(K_{Z_{\mathrm{sl}}}+cR_{\mathrm{sl}})|_F
                      R_{\mathrm{sl}}^{f_{\mathrm{sl}}-1},
 \end{split}
\end{equation}
with $c\le2(d+1-f_{\mathrm{sl}})/\delta$. Cartier multiples are
chosen before restriction, so the constant line from the fixed
factor causes no change in this calculation.

Let $W_0\subset X$ be the image of $F$ in the other factor.
There are three cases: $F$ is generically finite over a proper
subvariety of $X$; $F$ contains the projective-line fibres over its
image; or $F$ is a multisection over all of $X$.

\paragraph{A generically finite map to a proper subvariety.}
If $F\to W_0$ is generically finite and $\dim W_0<d$, use
$P_Y=R_{\mathrm{sl}}|_{F^*}$ and the pullback $L_Y$ from $X$.
The first inequality in \eqref{j:slice-bounds} gives a uniform
positive volume bound. Since $F$ meets neither axis generically,
its ambient canonical term is nonpositive. Hence the second
inequality gives a uniform canonical-intersection ratio.
Moreover $P_Y-rL_Y$ is nef. The image $W_0$ contains the very
general projection of the incidence point, so
\eqref{j:lower-big} applies. Lemma~\ref{lemma:bounded-curves}
excludes this case for large $r$.

\paragraph{A projective-line bundle over a proper subvariety.}
If $F\to W_0$ is not generically finite, its general fibre has
dimension one. Since the ambient projection is a projective-line
bundle and $F$ is integral, $F$ is the full bundle over $W_0$.
Write $w=\dim W_0=f_{\mathrm{sl}}-1$.
If $w=0$, then $F$ is a projective-line fibre and
$R_{\mathrm{sl}}\cdot F=q$. The first inequality in
\eqref{j:slice-bounds}, together with
\eqref{j:bundle-volumes} and \eqref{j:slice-gap}, gives
\[
 q\le(2/\delta)^nR_{\mathrm{sl}}^{d+1}
   \le(2/\delta)^d C_{\mathrm{sl},d}q\epsilon^d<q,
\]
which is impossible.

Suppose $w>0$. Resolve $W_0$ to $W_0^*$ and take as a resolution
of $F$ the induced projective bundle
$p_0:F^*\to W_0^*$. Its relative canonical divisor is
$-2\xi+p_0^*L$, exactly the restriction of
$K_{Z_{\mathrm{sl}}}$. Cancelling this term in
\eqref{j:slice-bounds} gives
\begin{equation}\label{j:saturated-canonical}
 p_0^*K_{W_0^*}R_{\mathrm{sl}}^w
                   \le cR_{\mathrm{sl}}^{w+1}\cdot F.
\end{equation}
The divisor $K_{W_0^*}$ is pseudo-effective by the very general
point condition. On expanding the left side and pushing down, its
term with exactly one $\xi$ is
$wqK_{W_0^*}P^{w-1}$; every other nonzero term is nonnegative,
since it pairs the pseudo-effective canonical divisor with nef
classes. On the right, the same expansion used for
\eqref{j:bundle-volumes} gives
\begin{equation}\label{j:saturated-volume}
 (w+1)qP^w\cdot W_0
 \le R_{\mathrm{sl}}^{w+1}\cdot F
 \le C'_nqP^w\cdot W_0,
\end{equation}
using $r\ge q$. The slice degree bound and the first inequality in
\eqref{j:saturated-volume} bound $0<P^w\cdot W_0$ uniformly.
Equations \eqref{j:saturated-canonical}--\eqref{j:saturated-volume}
also give
\[
 K_{W_0^*}P^{w-1}\le\frac{cC'_d}{w}P^w\cdot W_0.
\]
Apply \eqref{j:lower-big} and Lemma~\ref{lemma:bounded-curves}
on $W_0^*$ with the pullbacks of $P,L$. This excludes the second
case.

\paragraph{A multisection over $X$.}
The only remaining possibility is that $f_{\mathrm{sl}}=d$ and
$F\to X$ is generically finite of degree $e>0$. Intersect $F$ with
each of the two Cartier axes and push the resulting cycles down
to $X$. This gives effective integral Weil divisors $D_1,D_2$.
The axes do not contain $F$, and their classes differ by $L$;
projection of the associated rational section, or equivalently its
norm, gives
\begin{equation}\label{j:signed-L}
                         D_1-D_2\sim\pm eL.
\end{equation}
For either axis $S$, the classes $R_{\mathrm{sl}}-qS$ and
$R_{\mathrm{sl}}-P$ are nef. Intersecting with the effective cycle
$S\cdot F$ and then using nefness on $F$ yields
\begin{equation}\label{j:axis-degrees}
 qD_iP^{d-1}\le qS\cdot F\cdot R_{\mathrm{sl}}^{d-1}
                  \le R_{\mathrm{sl}}^d\cdot F.
\end{equation}
Thus both axis-image degrees are uniformly bounded.

Set $G=D_1+D_2$. It is $\Q$-Cartier because $X$ is
$\Q$-factorial. Choose $b>0$ rational sufficiently small that
$(X,bG)$ is klt. We claim
\begin{equation}\label{j:positive-boundary-kappa}
                         \kappa(K_X+bG)>0.
\end{equation}
The Iitaka dimension is at least zero, since
$K_X+bG\sim_\Q bG$. If it were zero, the model-program and ordinary abundance inputs in
$(\mathsf I_d)$ would give a finite ordinary MMP to a
good model $Y$ with transformed adjoint rationally linearly
trivial. The transform of $K_X$ is linearly trivial as well.
Consequently the effective transform of $G$ is zero: all its
components have been contracted. On a common resolution
$a:U\to X$, $b_Y:U\to Y$, relation \eqref{j:signed-L} then
expresses $a^*L$ up to rational linear equivalence as a signed
$b_Y$-exceptional divisor $E$. Divisors exceptional over $X$ are
also exceptional over $Y$, since the MMP extracts none.
Nefness of $a^*L$ and the negativity lemma imply $E\le0$.
Its class is also pseudo-effective. A nonzero anti-effective
divisor has negative intersection with a sufficiently high power
of an ample divisor, so $E=0$. This contradicts full nef dimension
of $L$, and proves \eqref{j:positive-boundary-kappa}.

The already proved positive-Iitaka-dimension case
(Proposition~\ref{prop:full-positive}) now gives
\[
                         K_X+bG+L\ \text{big}.
\]
Its full-positivity hypothesis holds because $L$ has full nef
dimension and $K_X+bG$ is pseudo-effective.
Take a log resolution $a:Y\to X$ of $G$ and let $D_0$ be the
reduced simple normal crossings divisor consisting of the strict
support of $G$ and every exceptional divisor. Decrease $b$ if
necessary; Proposition~\ref{prop:full-positive} still applies.
Terminality and the support of $a^*G$ then give
\[
 K_Y+D_0+a^*L\ \ge\ a^*(K_X+bG+L)
\]
up to rational linear equivalence and an effective difference.
Indeed the finitely many exceptional discrepancies are positive,
and decreasing $b$ makes each coefficient of $b\,a^*G$ no larger
than the corresponding coefficient of $K_Y-a^*K_X+D_0$.
Hence the left side is big.

Test on $Y$ with $a^*P$. Exceptional divisors contribute zero to
intersection with $(a^*P)^{d-1}$, and $K_X\sim0$. The degree of
the reduced strict support is no larger than $GP^{d-1}$, bounded
by \eqref{j:axis-degrees}; also
$LP^{d-1}\le P^d/r$. Since \eqref{j:volume-window} bounds $P^d$
away from zero, we obtain a uniform constant $C$ such that
\[
 (K_Y+D_0+a^*L)(a^*P)^{d-1}\le C(a^*P)^d.
\]
Lemma~\ref{lemma:bounded-curves}, now with its reduced boundary
$D_0$, gives the final contradiction. This proves
\eqref{j:no-intermediate-fibre}.

Thus a general moving component $V$ has fibre dimension either
zero or $d+1$ over its image in either factor. Fibre dimension
$d+1$ would make $V$ the full inverse image of its image in that
factor. Since $V$ is proper, the image has dimension at most
$d-1$, so the other projection would have fibre dimension between
$1$ and $d$, already excluded. Both projections of $V$ are
therefore generically finite onto their images, and $\dim V\le d$.
If $\dim V<d$, the bounds \eqref{j:total-component-bounds}, the
nef class $R-rL_i$, and \eqref{j:lower-big} again contradict
Lemma~\ref{lemma:bounded-curves}. Only $d$-dimensional
correspondences, dominant over both factors, remain.

\subsection{Correspondences and their branch divisors}

For these remaining $V$, \eqref{j:total-component-bounds} gives
uniform bounds
\begin{equation}\label{j:correspondence-bounds}
 R^d\cdot V\le C,
 \qquad K_{V^*}R^{d-1}\le cR^d\cdot V\le C,
\end{equation}
after increasing $C$. Each projection has bounded degree:
$R-P_i$ is nef, and $P^d$ is bounded below.

For each fixed $t$, resolve the general members in their algebraic
family and shrink the parameter space so that the resolutions form
a smooth projective family with integral fibres and morphisms to
$X$ in both slots. The dominant incidence and all numerical bounds
persist. Finite extensions of the parameter space are allowed.
Over $X_{\mathrm{sm}}$, name the components of the relative
Jacobian divisors as in Lemma~\ref{j:finite-covers}. Consider a
component whose image on a general member is a divisor $J\subset X$.

Such divisors $J$ cannot sweep $X$. If they did, choose a general
one through a very general point. Let $E_J$ be a ramification
divisor over it on the corresponding $V^*$. Pullback of a
canonical form from $X$ gives an effective canonical divisor on
$V^*$. Regularity follows by factoring through a resolution of
$X$ and using its canonical singularities; the coefficient of
$E_J$ is positive because its image meets $X_{\mathrm{sm}}$
generically. The difference $R-P_i$ is nef, so
\eqref{j:correspondence-bounds} and projection formula bound
\[
                    0<e_J:=P^{d-1}\cdot J\le C.
\]
On a resolution $J^*$, Lemma~\ref{j:subadjunction}, applied to
$(X,J)$ at the generic point of $J$, gives
\begin{equation}\label{j:branch-canonical}
 \begin{split}
 K_{J^*}P^{d-2}
 &\le(K_X+J)\cdot J\cdot P^{d-2}\\
 &=J^2P^{d-2}
 \le\frac{(JP^{d-1})^2}{P^d}
 \le C' e_J.
 \end{split}
\end{equation}
The pair is lc at that generic point because $X$ is normal and
hence smooth there. The middle inequality is the Hodge index
inequality, applied on a resolution and obtained by ample
approximation of the pulled-back testing class. It does not require
$J$ to be nef. The final constant is uniform by the volume lower
bound and the bound for $e_J$.

Since $J$ sweeps $X$, its resolution satisfies the very general
point conditions used in \eqref{j:lower-big}. With testing class
$P|_{J^*}$, whose top self-intersection is $e_J$, Equation
\eqref{j:branch-canonical} gives precisely the hypotheses of
Lemma~\ref{lemma:bounded-curves} in dimension $d-1$. This
contradiction excludes sweeping branch divisors.

Lemma~\ref{j:finite-covers} now gives, for each slot and this fixed
$t$, only finitely many possible finite normal Stein covers of
$X$, up to isomorphism over $X$. The point of the preceding step
is that degree bounds alone would not give this conclusion:
non-sweeping branch images provide a fixed smooth branch complement
for the family. Neither that complement nor the resulting finite
lists need be uniform in $t$.

\subsection{Fixed covers and an abelian model}

Each $V^*$ maps birationally to both of its normal Stein covers,
and hence determines a graph of a birational map between a pair
of covers in the finite lists. We explain why some such fixed pair
has a family of graphs dominating its product. Restrict the finite
set of pairs to those realized by at least one original member
$V^*$ satisfying \eqref{j:correspondence-bounds}. The original
incidence dominates $Z$, so the projected incidences for these
pairs cover a dense open of $X\times X$. For a fixed pair, graphs of
birational maps lie in countably many Hilbert schemes. Stratify
these schemes so that the fibres are integral and flat and both
projections are birational; the latter conditions may be tested
by dominance and degree one after shrinking. Because $\C$ is
uncountable, a dense open of $X\times X$ cannot be covered by
countably many proper closed subsets. Thus, for one fixed pair
$T_1,T_2$, some integral graph family has dominant evaluation to
$X\times X$. Since $T_1\times T_2$ is finite over $X\times X$
and has the same dimension, this evaluation is dominant onto
$T_1\times T_2$ as well.

Both covers are non-uniruled because they are finite over $X$.
The birational-group theorem of Hanamura
\cite[Theorems~2.1--2.2]{Hanamura}, in the form recalled in
\cite[Proposition~3.7 and Theorems~3.8--3.9]{Blanc}, implies that
they are birational to an abelian variety. Here are the details
of the application. Identify the two covers birationally by one
map in the family. On a suitable smooth projective birational
model of the resulting non-uniruled variety, the reduced
birational group is a group scheme locally of finite type, with
identity component an abelian variety acting regularly. Conjugate
the graph family to this model and shrink again to obtain a flat
family of integral graphs. By the flat-graph representability
statement, its reduced integral parameter space maps to the
reduced birational group. Its connected image lies in one
translate of the identity component. Product-dominant evaluation
therefore gives a dense orbit for that abelian identity component.
The orbit is a quotient of an abelian variety by a stabilizer,
hence is itself an abelian variety. The fixed cover is birational
to it. This is the birational-group part of the fixed-cover
argument in \cite[Proposition~6.11]{LA}; no canonical
nonvanishing conclusion from that proposition is being used.

This structural conclusion depends only on the two covers. We may
therefore return to an original member $V^*$ belonging to this pair,
which retains \eqref{j:correspondence-bounds}, and resolve
further to obtain a birational morphism
\[
                         p_0:V^*\longrightarrow A_0
\]
to an abelian variety. Higher resolutions do not change the
canonical intersection with the pulled-back $(d-1)$-st power of
$R$. Let $L_0$ be the pullback of $L$ from either slot. It is nef,
Cartier, and of full nef dimension. We claim
\begin{equation}\label{j:abelian-big}
                         K_{V^*}+L_0\ \text{big}.
\end{equation}

The pushforward $N_0=(p_0)_*L_0$ is pseudo-effective on $A_0$,
and is therefore nef. Indeed translations make the
pseudo-effective and nef cones of an abelian variety coincide.
If $N_0$ were not big, the standard description of nef line
bundles on abelian varieties would give a positive-dimensional
abelian subvariety on whose translates $N_0$ is numerically
trivial. One can see this from the semipositive Hermitian form of
the numerical class: its kernel is rational for the lattice,
since the alternating form of the integral Cartier class is
integral, and therefore defines an abelian subvariety.
A general translate meets the image of the exceptional locus of
$p_0$ in codimension at least two, if at all. Complete-intersection
curves in that translate can thus be chosen through very general
points and avoiding this image. Their lifts have $L_0$-degree
zero, contradicting full nef dimension. Hence $N_0$ is big.

By the negativity lemma,
\[
                         p_0^*N_0-L_0=E_0\ge0
\]
is exceptional. The canonical divisor of $V^*$ has an effective
representative $E_K$ with positive coefficient along every
$p_0$-exceptional divisor, because the target is smooth with
trivial canonical bundle. Choose a rational $0<u<1$ sufficiently
small that $E_K-uE_0\ge0$. Then
\[
 K_{V^*}+uL_0\sim_\Q u p_0^*N_0+(E_K-uE_0)
\]
is big. Adding the nef divisor $(1-u)L_0$ proves
\eqref{j:abelian-big}.

We may now apply Lemma~\ref{lemma:bounded-curves} in dimension
$d$ on $V^*$, with testing class $R$ and empty boundary.
The bounds \eqref{j:correspondence-bounds} give both the top
intersection bound and the canonical ratio, while $R-rL_0$ is
nef and \eqref{j:abelian-big} supplies bigness. This final
contradiction excludes the correspondence case.

\paragraph{Conclusion of the proof of
Proposition~\ref{prop:large-seshadri}.}
The alternatives for every moving base component have been
exhausted: first its slice fibres, then components of dimension
less than $d$, sweeping branch divisors, and finally the fixed
covers. In each case the contradiction comes from
Lemma~\ref{lemma:bounded-curves} with constants uniform in $t$.
The finite cover lists were used only for one fixed $t$, and their
sizes do not enter the degree bounds. Thus failures of
\eqref{j:large-seshadri} cannot occur along a sequence $t\to0$.
This proves the proposition, including the additional open
conditions on the chosen point.\hfill$\square$


\section{Proof of Lemma~\ref{lem:general-jets}}
\begin{proof}
Fix $s=q$ supplied by Proposition~\ref{prop:large-seshadri}.
For every $n\gg0$, choose $z\in U$ such that
\[
 \varepsilon(R_{n,s};z)>2d+2.
\]
Choose a rational number $a$ strictly between these two numbers.
On the blowup $b:\widetilde Z\to Z$ at $z$, with exceptional
divisor $G$, the divisor $b^*R_{n,s}-aG$ is ample. Serre vanishing
and the restriction sequence imply that, for all sufficiently
large divisible $k$, $|kR_{n,s}|$ separates $(ka-1)$-jets at $z$.
For such $k$, $ka-1\ge(2d+2)k$.

For one such $k_0$, this jet surjectivity is an open rank condition.
It therefore holds on a nonempty open subset of $U$. Products of
the degree-$k_0$ sections separate $\ell(2d+2)k_0$-jets in degree
$\ell k_0$ for every $\ell\ge1$. Indeed, every monomial of total
degree at most $\ell(2d+2)k_0$ is a product of $\ell$ monomials
of total degree at most $(2d+2)k_0$, and the lifted products span
the jet quotient. This proves the assertion at a very general
point and for all sufficiently divisible $k$.
\end{proof}

\begin{thebibliography}{99}
\bibitem{NS}
\emph{Numerical Semiampleness of Nef Adjoint Divisors},
OpenAI mathematics collection, October 3, 2026,
Sections~2--6, especially Proposition~5.3 and Theorem~6.1.
\href{https://github.com/openai/math/blob/main/preprints/Numerical-Semiampleness-of-Nef-Adjoint-Divisors-October-3-2026/paper.pdf}{Source manuscript}.

\bibitem{BZ}
C.~Birkar and D.-Q.~Zhang,
\emph{Effectivity of Iitaka fibrations and pluricanonical systems of polarized pairs},
Publ. Math. Inst. Hautes \'Etudes Sci. \textbf{123} (2016), 283--331.
\href{https://doi.org/10.1007/s10240-016-0080-x}{DOI}.

\bibitem{Blanc}
J.~Blanc,
\emph{Algebraic structures of groups of birational transformations},
in \emph{Algebraic Groups: Structure and Actions} (M.~Bilen Can, ed.),
Proc.~Sympos.~Pure Math. \textbf{94}, American Mathematical Society,
Providence, RI, 2017, pp.~17--30.
\href{https://algebra.dmi.unibas.ch/blanc/articles/structurebir.pdf}{Author-hosted published offprint}.

\bibitem{BDPP}
S.~Boucksom, J.-P.~Demailly, M.~P\u{a}un, and Th.~Peternell,
\emph{The pseudo-effective cone of a compact K\"ahler manifold and varieties of negative Kodaira dimension},
J.~Algebraic Geom. \textbf{22} (2013), no.~2, 201--248.
\href{https://doi.org/10.1090/S1056-3911-2012-00574-8}{DOI}.

\bibitem{FG}
O.~Fujino and Y.~Gongyo,
\emph{On the moduli $b$-divisors of lc-trivial fibrations},
Ann. Inst. Fourier (Grenoble) \textbf{64} (2014), no.~4, 1721--1735.
\href{https://doi.org/10.5802/aif.2894}{DOI}.

\bibitem{FH}
O.~Fujino and K.~Hashizume,
\emph{Existence of log canonical modifications and its applications},
Eur. J. Math. \textbf{9} (2023), no.~1, Paper No.~13, 32~pp.
\href{https://doi.org/10.1007/s40879-023-00598-0}{DOI}.

\bibitem{Hanamura}
M.~Hanamura,
\emph{Structure of birational automorphism groups, I: non-uniruled varieties},
Invent.~Math. \textbf{93} (1988), no.~2, 383--403.
\href{https://doi.org/10.1007/BF01394338}{DOI}.

\bibitem{LazarsfeldPAGI}
R.~Lazarsfeld,
\emph{Positivity in Algebraic Geometry I: Classical Setting: Line Bundles and Linear Series},
Ergebnisse der Mathematik und ihrer Grenzgebiete, vol.~48,
Springer-Verlag, 2004.
\href{https://www.math.stonybrook.edu/~roblaz/Reprints/Laz.PAG.Chapt1.pdf}{Author-hosted Chapter~1}.

\bibitem{LA}
OpenAI,
\emph{Log abundance in characteristic zero},
OpenAI Math Release preprint
\href{https://github.com/openai/math/blob/main/preprints/Log-abundance-in-characteristic-zero-September-24-2026/paper.pdf}{OAI:Log-abundance-in-characteristic-zero-September-24-2026},
2026.
\end{thebibliography}

\end{document}
